\documentclass[10pt,a4paper]{article}
\usepackage[margin=20mm,headheight=15pt,headsep=7mm,footskip=10mm]{geometry}
\usepackage{fontspec,amsmath,amssymb,mathtools,microtype}
\setmainfont{lmroman10-regular.otf}[BoldFont=lmroman10-bold.otf,ItalicFont=lmroman10-italic.otf,BoldItalicFont=lmroman10-bolditalic.otf]
\setsansfont{lmsans10-regular.otf}[BoldFont=lmsans10-bold.otf,ItalicFont=lmsans10-oblique.otf,BoldItalicFont=lmsans10-boldoblique.otf]
\setmonofont[Scale=.86]{lmmono10-regular.otf}[ItalicFont=lmmono10-italic.otf]
\usepackage{ragged2e}\usepackage{xcolor,booktabs,longtable,array,graphicx,enumitem,fancyhdr,needspace,titlesec,hyperref,xurl}
\definecolor{AtlasNavy}{HTML}{18334C}\definecolor{AtlasTeal}{HTML}{116D73}\definecolor{AtlasGray}{HTML}{52616D}
\hypersetup{colorlinks=true,linkcolor=AtlasTeal,urlcolor=AtlasTeal,pdfcreator={Complexity Atlas},unicode=true}
\urlstyle{same}\setlength{\parindent}{0pt}\setlength{\parskip}{4.8pt}\setlength{\emergencystretch}{3em}
\renewcommand{\baselinestretch}{1.08}\setcounter{secnumdepth}{0}\setlist{nosep,leftmargin=1.5em}
\titleformat{\section}{\Large\bfseries\sffamily\color{AtlasNavy}\raggedright}{}{0em}{}
\titleformat{\subsection}{\large\bfseries\sffamily\color{AtlasTeal}\raggedright}{}{0em}{}
\titleformat{\subsubsection}{\normalsize\bfseries\color{AtlasNavy}\raggedright}{}{0em}{}
\titlespacing*{\section}{0pt}{14pt}{6pt}\titlespacing*{\subsection}{0pt}{10pt}{5pt}
\pagestyle{fancy}\fancyhf{}\fancyhead[L]{\small\sffamily\color{AtlasGray}COMPLEXITY ATLAS}
\fancyhead[R]{\small\sffamily\color{AtlasGray}ENGLISH EDITION}\fancyfoot[C]{\small\thepage}\renewcommand{\headrulewidth}{.3pt}
\newcommand{\AtlasMath}[1]{\par\begingroup\centering\setbox0=\hbox{$\displaystyle #1$}\ifdim\wd0>1.4\linewidth\typeout{ATLAS WIDE FORMULA: \the\wd0}\fi\ifdim\wd0>\linewidth\resizebox{\linewidth}{!}{\box0}\else\box0\fi\par\endgroup}
\newcommand{\AtlasInlineMath}[1]{\begingroup\setbox0=\hbox{$#1$}\ifdim\wd0>.4\linewidth\endgroup\AtlasMath{#1}\else\endgroup$#1$\fi}
\pdfstringdefDisableCommands{\def\AtlasInlineMath#1{#1}}
\DeclareMathOperator{\polylog}{polylog}\DeclareMathOperator{\dist}{dist}\DeclareMathOperator{\diam}{diam}
\DeclareMathOperator{\rank}{rank}\DeclareMathOperator{\tr}{tr}\DeclareMathOperator{\diag}{diag}
\DeclareMathOperator{\prox}{prox}\DeclareMathOperator{\KL}{KL}\DeclareMathOperator{\Var}{Var}
\providecommand{\eps}{\varepsilon}\providecommand{\R}{\mathbb R}\providecommand{\E}{\mathbb E}
\providecommand{\mathbbm}[1]{\mathbf{#1}}\providecommand{\mathds}[1]{\mathbf{#1}}
\providecommand{\llbracket}{[\![}\providecommand{\rrbracket}{]\!]}
\begin{document}
\raggedright

{\LARGE\bfseries\sffamily\color{AtlasNavy} General nonconvex-nonconcave optimization: exact local minimax points.\par}\vspace{5pt}
\textbf{Setting ID.} \texttt{CV\_\allowbreak{}S\_\allowbreak{}LOCAL\_\allowbreak{}MINIMAX}

Minimax optimization

Find an exact local minimax point \AtlasInlineMath{(x^\star,y^\star)\in\mathbb R^{d_x}\times\mathbb R^{d_y}}. There must exist \AtlasInlineMath{r_0>0} and a function \AtlasInlineMath{h:(0,r_0]\to[0,\infty)} with \AtlasInlineMath{\lim_{r\downarrow0}h(r)=0} such that every \AtlasInlineMath{0<r\le r_0}, \AtlasInlineMath{\|x-x^\star\|_2\le r}, and \AtlasInlineMath{\|y-y^\star\|_2\le r} satisfies

\textbf{Status.} Infeasible task

\textbf{TASK tags.} formulation: Local sequential minimax; solution concept: Exact local minimax; task family: Minimax

\textbf{MODEL tags.} convexity: C2 nonconvex nonconcave; existence: Existence not promised

\textbf{ORACLE tags.} feedback quality: Exact; interface family: Value joint gradient; operation: Payoff value joint gradient; public information: Hidden objects explicitly declared; public information: Public inputs explicitly declared; query domain: Full space; scalar derivative order: Scalar-function first derivative

\textbf{ALGORITHMS tags.} adaptivity: Adaptive; randomization: Deterministic; restriction: Explicitly no additional algorithm restrictions; restriction: No additional span restriction explicitly imposed

\textbf{OUTPUT tags.} certificate: No additional output certificate required; object: Primal dual pair; postprocessing: Measurable postprocessing; postprocessing: Unqueried output explicitly allowed

\textbf{METRIC tags.} criterion: Local minimax exact

\textbf{GUARANTEE tags.} quantifier: Algorithm before instance

\textbf{RESOURCE tags.} budget: Fixed worst-case budget; measure: N bundle; precision: Bit precision not counted

\needspace{7\baselineskip}\section{1. Setting specification}
\needspace{4\baselineskip}\subsection{Task}
Find an exact local minimax point \AtlasInlineMath{(x^\star,y^\star)\in\mathbb R^{d_x}\times\mathbb R^{d_y}}. There must exist \AtlasInlineMath{r_0>0} and a function \AtlasInlineMath{h:(0,r_0]\to[0,\infty)} with \AtlasInlineMath{\lim_{r\downarrow0}h(r)=0} such that every \AtlasInlineMath{0<r\le r_0}, \AtlasInlineMath{\|x-x^\star\|_2\le r}, and \AtlasInlineMath{\|y-y^\star\|_2\le r} satisfies


\AtlasMath{f(x^\star,y)\le f(x^\star,y^\star)\le\max_{\|y^\prime-y^\star\|_2\le h(r)} f(x,y^\prime).}


The two neighborhood radii are r and h(r). Neither h(r)=r nor proportionality to r is required, and the radii are not supplied in advance.

\needspace{4\baselineskip}\subsection{Admissible model}
Assume \AtlasInlineMath{f\in C^2(\mathbb{R}^{d_x+d_y})}; neither variable block must be convex or concave. Existence of a local minimax point is not promised, and no quantitative regularity, neighborhood radius, or approximate-solution tolerance is specified.

\needspace{4\baselineskip}\subsection{Oracle interface}
Public inputs are the positive integer block dimensions \AtlasInlineMath{d_x,d_y}. The analytic expression for \AtlasInlineMath{f} is hidden. The output must be exact; there is no adjustable positive error tolerance.


\AtlasMath{(x,y)\longmapsto(f(x,y),\nabla_x f(x,y),\nabla_y f(x,y)).}


Exact full-space queries.

\needspace{4\baselineskip}\subsection{Output}
Return an exact local minimax point \AtlasInlineMath{(\widehat x,\widehat y)\in\mathbb R^{d_x+d_y}}.

\needspace{4\baselineskip}\subsection{Loss or accuracy criterion}
The set \AtlasInlineMath{\mathcal L_f} contains every local minimax point under the stated local definition.


\AtlasMath{\mathbf 1\{(\widehat x,\widehat y)\notin\mathcal L_f\}.}


The indicator loss must equal zero. Here \AtlasInlineMath{\varepsilon=0} denotes an exact task, not a gradient-norm tolerance.

\needspace{4\baselineskip}\subsection{Resource accounting}
Count bundled value-and-joint-gradient calls. Because existence is not guaranteed, no uniform finite-query upper bound results.

Count all information calls for one instance and one requested accuracy, including initialization. Exact-real arithmetic, storage, and solving known models from public inputs and returned data are uncharged. This query-complexity quantity does not certify arithmetic or bit complexity.

\needspace{4\baselineskip}\subsection{Algorithm views and success guarantees}
Any measurable sequential adaptive deterministic algorithm is allowed. Queries and outputs may use the complete history, with no additional span, zero-support, communication, or memory restriction.

A single deterministic algorithm must return a point in \AtlasInlineMath{\mathcal L_f} for every admissible \AtlasInlineMath{f}. The model allows \AtlasInlineMath{\mathcal L_f=\varnothing}, so this uniform requirement is infeasible.

\textbf{Budget convention.} fixed worst-case query cap

\needspace{7\baselineskip}\section{2. Notation and parameters}
\AtlasInlineMath{d}: The joint Euclidean dimension is \AtlasInlineMath{d=d_x+d_y}. Within its stated window, the lower bound may choose finite block dimensions. \textbf{Domain:} Integer at least \AtlasInlineMath{2}.

\AtlasInlineMath{\varepsilon}: Fixed zero tolerance for exact decision \textbf{Domain:} Fixed at zero for exact decision.

\AtlasInlineMath{\varepsilon=0}

\AtlasInlineMath{\|\cdot\|_2,\langle\cdot,\cdot\rangle}: Euclidean norm and inner product; the product-space norm is the square root of the sum of the squared block norms. \textbf{Domain:} The specified finite-dimensional real vector space.

\AtlasInlineMath{\|v\|_2=\sqrt{\langle v,v\rangle}}

\AtlasInlineMath{d_x,d_y}: Positive integer primal and dual block dimensions; joint dimension \AtlasInlineMath{d=d_x+d_y} \textbf{Domain:} Positive integer.

\AtlasInlineMath{f}: Unknown real-valued payoff, with primal and dual input blocks. \textbf{Domain:} \AtlasInlineMath{f:(x,y)\mapsto f(x,y)\in\mathbb R}.

\AtlasInlineMath{x,y,z}: Primal variable, dual variable, and joint point \textbf{Domain:} \AtlasInlineMath{x\in\mathbb R^{d_x},\ y\in\mathbb R^{d_y},\ z=(x,y)\in\mathbb R^{d_x+d_y}}.

\AtlasInlineMath{z=(x,y)}

\AtlasInlineMath{\nabla f,\nabla_x f,\nabla_y f}: Complete joint gradient and its two blocks \textbf{Domain:} \AtlasInlineMath{\nabla_x f\in\mathbb R^{d_x},\ \nabla_y f\in\mathbb R^{d_y}}.

\AtlasInlineMath{\nabla f=(\nabla_x f,\nabla_y f)}

\AtlasInlineMath{\widehat x,\widehat y}: The point or variable blocks ultimately returned by the algorithm, corresponding to the hatted variables in the error formula. \textbf{Domain:} Feasible variable pairs in \AtlasInlineMath{\mathbb R^{d_x}\times\mathbb R^{d_y}}.

\AtlasInlineMath{N_{\rm first}}: Number of bundled function-value and full joint-gradient queries \textbf{Domain:} Nonnegative integer count; expected work is its expectation.

\AtlasInlineMath{N}: This is the principal information-call budget on one instance. The resource specification determines whether it is fixed or a worst-instance expected budget. \textbf{Domain:} A nonnegative integer; expected budgets may be nonnegative real numbers.

\AtlasInlineMath{N^\star}: The infimum principal query budget among algorithms satisfying every stated instancewise output guarantee. \textbf{Domain:} Nonnegative extended real value; by convention, positive infinity if no successful algorithm or no uniform finite budget exists.

\AtlasInlineMath{\mathcal A,I}: An algorithm selected from public inputs and an admissible hidden instance fixed before interaction \textbf{Domain:} The specified algorithm and model classes.

\AtlasInlineMath{\mathcal L_f}: The exact local minimax points under the complete local definition stated here. \textbf{Domain:} A subset of \AtlasInlineMath{\mathbb R^{d_x+d_y}}, possibly empty.

\AtlasInlineMath{x^\star,y^\star}: An arbitrary candidate point in the local definition, not an assumed global solution. \textbf{Domain:} \AtlasInlineMath{\mathbb R^{d_x}} and \AtlasInlineMath{\mathbb R^{d_y}}.

\AtlasInlineMath{r_0,r,h}: The local threshold, vanishing neighborhood radius, and dual-response radius function. \textbf{Domain:} \AtlasInlineMath{r_0>0,\ 0<r\le r_0,\ h:(0,r_0]\to[0,\infty)}.

\AtlasInlineMath{\lim_{r\downarrow0}h(r)=0}

\AtlasInlineMath{\mathbf1\{E\}}: The indicator of proposition E: one if true and zero otherwise. \textbf{Domain:} \AtlasInlineMath{\{0,1\}}.

\AtlasInlineMath{e_1}: First coordinate unit vector in the primal-variable space \textbf{Domain:} \AtlasInlineMath{\mathbb R^{d_x}}.

\AtlasInlineMath{\mathfrak A_{\rm succ}}: Algorithms returning an exact local minimax point for every admissible function \textbf{Domain:} Empty set for the task over the full model class.

\AtlasInlineMath{N_{\mathcal A}(I,\omega)}: Primary information calls accumulated by algorithm A on execution path omega of fixed instance I \textbf{Domain:} A nonnegative integer. For an almost surely terminating algorithm, use its complete stopping count.

\needspace{12\baselineskip}\section{3. Complexity bracket}
\needspace{4\baselineskip}\subsection{Compared quantities}
\needspace{4\baselineskip}\subsubsection{Q1 - target accuracy cost}
\AtlasMath{N^\star=\inf_{\mathcal A\in\mathfrak A_{\rm succ}}\sup_I N_{\mathcal A}(I),\qquad\inf\varnothing:=+\infty}
\textbf{Scope.} A successful algorithm must return an exact local minimax point for every admissible function. A linear counterexample makes the admissible algorithm set empty; this is not an ordinary finite-error query rate.

\textbf{Optimization direction.} minimize

\needspace{4\baselineskip}\subsection{R1 - Full model class, which contains instances with no solution}
\textbf{Region.} All declared \AtlasInlineMath{C^2} nonconvex-nonconcave full-space functions, without a solution-existence promise.

\textbf{Quantity.} Q1

\textbf{Resource.} The principal resource specified in the resource section. Primitive operations and unit costs are measured along this axis.

\textbf{Status.} Infeasible task

\textbf{Lower bound.} B1

\AtlasMath{N^\star=+\infty}
\textbf{Comparison argument.} No algorithm succeeds on the entire class because some admissible instances have no solution. This is task infeasibility, not missing upper and lower bounds.

\needspace{7\baselineskip}\section{4. Bounds and supporting evidence}
The following conditions apply to every bound in this profile, in addition to each bound's stated region and algorithm scope.

\textbf{Conditions.} For arbitrary positive integers \AtlasInlineMath{d_x,d_y}, let \AtlasInlineMath{f(x,y)=x_1}. This function is \AtlasInlineMath{C^\infty} on the whole space and satisfies the stated \AtlasInlineMath{C^2} assumptions. The model imposes no lower boundedness, compact primal domain, strong concavity, or existence promise. For any candidate \AtlasInlineMath{(x^\star,y^\star)}, any \AtlasInlineMath{r_0>0}, and any admissible \AtlasInlineMath{h}, choose \AtlasInlineMath{0<r\le r_0} and \AtlasInlineMath{x=x^\star-(r/2)e_1}. For every \AtlasInlineMath{y^\prime}, \AtlasInlineMath{f(x,y^\prime)=x_1^\star-r/2<f(x^\star,y^\star)}, violating the second inequality in the local minimax definition. Thus an instance without a solution exists in every allowed dimension; even revealing the entire function cannot produce a valid output. The notation \AtlasInlineMath{N^\star=+\infty} uses the convention that the successful-algorithm class is empty. The explanatory status is \texttt{infeasible\_\allowbreak{}task}.

\textbf{Dimension conditions.} Any \AtlasInlineMath{d_x,d_y\ge1} are allowed. The counterexample already holds at \AtlasInlineMath{d_x=d_y=1}.

\textbf{Oracle mapping.} The counterexample is independent of the information available. Exact-value and gradient interfaces satisfy the model, and additional information cannot create a nonexistent point.

\textbf{Guarantee mapping.} The requirement that one algorithm succeed on every instance fails because the class contains an admissible instance with no solution.

\textbf{Resource mapping.} No algorithm succeeds, so neither a finite nor an infinite query budget satisfies the task on the entire class.

\needspace{4\baselineskip}\subsection{B1 - No-solution counterexample with a linear primal objective}
\textbf{Admission.} Included in the bracket

\textbf{Direction.} lower

\textbf{Witness or lower-bound scope.} All randomized adaptive algorithms

\textbf{Algorithm-scope qualification.} An admissible instance has an empty solution set, so no output strategy can succeed. No span, zero-respecting, or queried-output restriction is needed.

\AtlasMath{N^\star=+\infty}
\textbf{Sources.} \href{https://proceedings.mlr.press/v119/jin20e/jin20e.pdf}{[1]}

\textbf{Constants.} This no-solution counterexample involves no asymptotic constant or accuracy threshold.

\textbf{Local argument.} A1

\textbf{Source location.} \href{https://proceedings.mlr.press/v119/jin20e/jin20e.pdf}{[1]} theorem: Definition 14; section: Section 3.1; pdf pages: 5

Use the original definition’s quantifiers. The source paper supports the definition only; the counterexample f = x\_1 is a separate argument, not a theorem attributed to that paper.

\needspace{7\baselineskip}\section{5. Local arguments}
\needspace{4\baselineskip}\subsection{A1 - No-solution counterexample with a linear primal objective}
\textbf{Scope.} All model, oracle, algorithm, guarantee, quantity, resource-projection, and region conditions attached to the mapped bounds remain in force.

For arbitrary positive integers \AtlasInlineMath{d_x,d_y}, let \AtlasInlineMath{f(x,y)=x_1}. This function is \AtlasInlineMath{C^\infty} on the whole space and satisfies the stated \AtlasInlineMath{C^2} assumptions. The model imposes no lower boundedness, compact primal domain, strong concavity, or existence promise. For any candidate \AtlasInlineMath{(x^\star,y^\star)}, any \AtlasInlineMath{r_0>0}, and any admissible \AtlasInlineMath{h}, choose \AtlasInlineMath{0<r\le r_0} and \AtlasInlineMath{x=x^\star-(r/2)e_1}. For every \AtlasInlineMath{y^\prime}, \AtlasInlineMath{f(x,y^\prime)=x_1^\star-r/2<f(x^\star,y^\star)}, violating the second inequality in the local minimax definition. Thus an instance without a solution exists in every allowed dimension; even revealing the entire function cannot produce a valid output. The notation \AtlasInlineMath{N^\star=+\infty} uses the convention that the successful-algorithm class is empty. The explanatory status is \texttt{infeasible\_\allowbreak{}task}.

\textbf{Sources.} \href{https://proceedings.mlr.press/v119/jin20e/jin20e.pdf}{[1]}

\needspace{7\baselineskip}\section{References}
[1] Chi Jin; Praneeth Netrapalli; Michael I. Jordan. What is Local Optimality in Nonconvex-Nonconcave Minimax Optimization?. 2020. ICML PMLR119.

\url{https://proceedings.mlr.press/v119/jin20e/jin20e.pdf}

\end{document}
